\documentclass[../../../main/main.tex]{subfiles}

\begin{document}

\chapter{Linear Notions}

\begin{theorem}[Infinite interior points]
  \label{lem:infinite-interior}
  Let $A,C\in\mathbf{P}$ with $A\neq C$.  Then there exist infinitely many
  points strictly between $A$ and $C$.
\end{theorem}
\begin{proof}
    wasd
\end{proof}

\begin{theorem}[Infinite outward extension]
  \label{lem:infinite-outward}
  Let $A,B\in\mathbf{P}$ with $A\neq B$.  Then there exist infinitely many
  points $C_1,C_2,C_3,\dots$ such that
  \[
  \mathcal{B}(A,B,C_1),\quad
  \mathcal{B}(A,C_1,C_2),\quad
  \mathcal{B}(A,C_2,C_3),\quad\dots
  \]
\end{theorem}
\begin{proof}
  wasd
\end{proof}

Together, these two lemmas show that from any segment $AB$ ($A\neq B$) we
can generate infinitely many points in both directions along the line
determined by $A$ and $B$.  Thus the plane $\mathcal{P}$ is necessarily
infinite.

\section{Defintion of Linear Notions}
A segment is a subset of $\mathbf{P}$ defined as the union of two discint points and all points between them. Symbolically, that is
\begin{definition}[Segment]
    \[
    \forall A, B \in \mathbf{P}, A \neq B \implies AB \coloneqq \{A, B\} \cup \{X \in \mathbf{P} \mid \mathcal{B}(A, X, B)\}.
    \]
\end{definition}
A ray pointing from point $A$ to a discint point $B$ is a subset of $\mathbf{P}$ denoted as $\overrightarrow{AB}$. It is defined as the union of $A$, $B$, all points between $A$ and $B$, and all points such that $B$ is between $A$ and that point, that is
\begin{definition}[Ray]
    \[
    \forall A, B \in \mathbf{P}, A \neq B \implies \overrightarrow{AB} \coloneqq \{A, B\} \cup \{X \in \mathbf{P} \mid \mathcal{B}(A, X, B) \lor \mathcal{B}(A, B, X)\}.
    \]
\end{definition}
The line $\overleftrightarrow{AB}$ is the subset of $\mathbf{P}$ consisting of the points $A$, $B$, all points between $A$ and $B$, all points where $A$ is between it and $B$, and all points where $B$ is between it and $A$. That is,
\begin{definition}[Line]
    \[
    \forall A, B \in \mathbf{P}, A \neq B \implies \overleftrightarrow{AB} \coloneqq \{A, B\} \cup \{X \in \mathbf{P} \mid \mathcal{B}(X, A, B) \lor \mathcal{B}(A, X, B) \lor \mathcal{B}(A, B, X)\}.
    \]
\end{definition}

\section{Properties of Linear Notions}

\begin{theorem}[Distinctness of Betweenness]
    If $\mathcal{B}(A, B, X)$ holds, then $A$, $B$, and $X$ are three distinct points:
    \[
    \forall A, B, X \in \mathbf{P}, \mathcal{B}(A, B, X) \implies (A \neq B \land B \neq X \land A \neq X).
    \]
\end{theorem}
\begin{proof}
    We proceed by contradiction. Suppose $\mathcal{B}(A, B, X)$ holds.

    First, assume $A = X$. Substituting $A$ for $X$ in our hypothesis, we obtain $\mathcal{B}(A, B, A)$, which  implies $A = B$. If $A = B$, then the relation $\mathcal{B}(A, B, X)$ is equivalent to $\mathcal{B}(B, A, X)$. However, by the Order of Betweeness, $\mathcal{B}(A, B, X) \implies \neg \mathcal{B}(B, A, X)$. Here, we reach a contradiction and conclude our hypothesis must be false and thus $A \neq X$.

    Next, we test the assumption $A = B$. Under this assumption, we substite our hypothesis to obtain $\mathcal{B}(B, A, X)$. Similarly, the Order of Betweeness shows these cannot be simutanously true, and thus our assumption must be false, meaning $A \neq B$.

    Finally, by the Transitivity of Equality, $B \neq X$. Thus proving the result.
\end{proof}

\begin{theorem}[Segment Symmetry]
    For any two points $A, B$, the segment $AB$ is equal to the segment $BA$.
    \[
    \forall A, B \in \mathbf{P}, AB = BA.
    \]
\end{theorem}
\begin{proof}
    By the definition of a line segment, we have:
    \[
    AB = \{A, B\} \cup \{X \in \mathbf{P} \mid \mathcal{B}(A, X, B)\}
    \]
    and
    \[
    BA = \{B, A\} \cup \{X \in \mathbf{P} \mid \mathcal{B}(B, X, A)\}.
    \]
    First, by the commutative property of sets, the endpoints are identical: $\{A, B\} = \{B, A\}$.

    Second, we examine the sets of interior points. By the \textbf{Symmetry of Betweenness}, the relation $\mathcal{B}(A, X, B)$ holds if and only if $\mathcal{B}(B, X, A)$ holds. Formally, this implies:
    \[
    \{X \in \mathbf{P} \mid \mathcal{B}(A, X, B)\} = \{X \in \mathbf{P} \mid \mathcal{B}(B, X, A)\}.
    \]
    Since $AB$ and $BA$ are defined as the union of identical sets, they contain the same points. By the principle of extensionality in set theory, $AB = BA$.
\end{proof}

\begin{theorem}[Line Symmetry]
    For any two distinct points $A$ and $B$, the line $\overleftrightarrow{AB}$ is equal to the segment $\overleftrightarrow{BA}$:
    \[
    \forall A, B \in \mathbf{P}, \overleftrightarrow{AB} = \overleftrightarrow{BA}.
    \]
\end{theorem}
\begin{proof}
    By the definition of a line, we have:
    \[
    \overleftrightarrow{AB} = \{A, B\} \cup \{X \in \mathbf{P} \mid \mathcal{B}(X, A, B) \lor \mathcal{B}(A, X, B) \lor \mathcal{B}(A, B, X)\}
    \]
    and
    \[
    \overleftrightarrow{BA} = \{B, A\} \cup \{X \in \mathbf{P} \mid \mathcal{B}(X, B, A) \lor \mathcal{B}(B, X, A) \lor \mathcal{B}(B, A, X)\}.
    \]
    First, by the commutative property of sets, the endpoints are identical: $\{A, B\} = \{B, A\}$.

    Second, we demonstrate that the three cases of betweeness in $\overleftrightarrow{AB}$ are logically equivalent to those of $\overleftrightarrow{BA}$ using the by the \textbf{Symmetry of Betweenness}:
    \begin{itemize}
        \item $\{X \in \mathbf{P} \mid \mathcal{B}(X, A, B)\} = \{X \in \mathbf{P} \mid \mathcal{B}(B, A, X)\}$
        \item $\{X \in \mathbf{P} \mid \mathcal{B}(A, X, B)\} = \{X \in \mathbf{P} \mid \mathcal{B}(B, X, A)\}$
        \item $\{X \in \mathbf{P} \mid \mathcal{B}(A, B, X)\} = \{X \in \mathbf{P} \mid \mathcal{B}(X, B, A)\}$
    \end{itemize}
    Since each condition for membership in $\overleftrightarrow{AB}$ is logically equivalent to a condition for membership in $\overleftrightarrow{BA}$, the sets contain the same points. Thus, by the principle of extensionality, $\overleftrightarrow{AB} = \overleftrightarrow{BA}$
\end{proof}

\begin{theorem}[Ray Asymmetry]
    For any two distinct points $A, B \in \mathbf{P}$, the ray $\overrightarrow{AB}$ is not equal to the ray $\overleftarrow{BA}$:
    \[
    \forall A, B \in \mathbf{P}, A \neq B \implies \overrightarrow{AB} \neq \overrightarrow{BA}.
    \]
\end{theorem}
\begin{proof}
    We proceed by by contradiction. Assume that the rays for some distinct points $A, B \in \mathbf{P}$, $\overrightarrow{AB} = \overrightarrow{BA}$.

    By the definition of a ray, we have:
    \[
    \overrightarrow{AB} = \{A, B\} \cup \{X \in \mathbf{P} \mid \mathcal{B}(A, X, B) \lor \mathcal{B}(A, B, X)\},
    \]
    and
    \[
    \overrightarrow{BA} = \{B, A\} \cup \{X \in \mathbf{P} \mid \mathcal{B}(B, X, A) \lor \mathcal{B}(B, A, X)\}.
    \]
    By the \textbf{Outer Point Axiom}, there must exist a point $Y$ such that $\mathcal{B}(A, B, Y)$. By the defintion of $\overrightarrow{AB}$, it follows that $Y \in \overrightarrow{AB}$. Under our initial assumption that the rays, $Y$ must be an element of $\overrightarrow{BA}$.

    By the Distinctness of Betweeness, the relation $\mathcal{B}(A, B, Y)$ implies that $A \neq Y$ and $B \neq Y$. Therefore, $Y \notin \{A, B\}$.

    Since $Y \in \overrightarrow{BA}$ but is not an endpoint, it must satisfy $\mathcal{B}(B, Y, A)$ or $\mathcal{B}(B, A, Y)$.

    However, because we know that $\mathcal{B}(A, B, Y)$ is true, by Order of Betweeness we have $\mathcal{B}(A, B, Y) \implies \neg \mathcal{B}(B, A, Y)$, and $\mathcal{B}(A, B, Y) \implies \neg \mathcal{B}(A, Y, B)$, which by the Symmetry of Betweeness is logically equivalent to $\mathcal{B}(A, B, Y) \implies \neg \mathcal{B}(B, Y, A)$.

    Thus, both conditions for $Y$ to be an element of $\overrightarrow{BA}$ are false. Therefore, the rays $\overrightarrow{AB}$ and $\overrightarrow{BA}$ are not equal, thus proving the desired result.
\end{proof}

\begin{theorem}[Union of Opposite Rays]
    For any two points $A, B \in \mathbf{P}$,
    \[
    \overleftarrow{AB} \cup \overrightarrow{AB} = \overleftrightarrow{AB}.
    \]
\end{theorem}
\begin{proof}
    By the definition of a ray, we have:
    \[
    \overleftarrow{AB} = \{A, B\} \cup \{X \in \mathbf{P} \mid \mathcal{B}(X, A, B) \lor \mathcal{B}(A, X, B)\},
    \]
    and
    \[
    \overrightarrow{AB} = \{A, B\} \cup \{X \in \mathbf{P} \mid \mathcal{B}(A, X, B) \lor \mathcal{B}(A, B, X)\}.
    \]
    The union of $\overleftarrow{AB}$ and $\overrightarrow{AB}$ therefore consists of the set, $\{B, A\}$, and all points $X$ such that $\mathcal{B}(X, A, B)$ or $\mathcal{B}(A, X, B)$, or $\mathcal{B}(A, B, X)$.

    Note that this is is the same as the line $\overleftrightarrow{AB}$, as it is defined as the set
    \[
    \overleftrightarrow{AB} = \{A, B\} \cup \{X \in \mathbf{P} \mid \mathcal{B}(X, A, B) \lor \mathcal{B}(A, X, B) \lor \mathcal{B}(A, B, X)\}.
    \]
    This set is equal to the union described above. Thus, the union of the rays $\overleftarrow{AB}$ and $\overrightarrow{AB}$ is equal to the line $\overleftrightarrow{AB}$.
\end{proof}

\begin{theorem}[Intersection of Opposite Rays]
    For any two points $A, B \in \mathbf{P}$,
    \[
    \overleftarrow{AB} \cap \overrightarrow{AB} = AB.
    \]
\end{theorem}
\begin{proof}
    By the definition of a ray, we have:
    \[
    \overleftarrow{AB} = \{A, B\} \cup \{X \in \mathbf{P} \mid \mathcal{B}(X, A, B) \lor \mathcal{B}(A, X, B)\},
    \]
    and
    \[
    \overrightarrow{AB} = \{A, B\} \cup \{X \in \mathbf{P} \mid \mathcal{B}(A, X, B) \lor \mathcal{B}(A, B, X)\}.
    \]
    The intersection of $\overleftarrow{AB}$ and $\overrightarrow{AB}$ therefore consists of the set, $\{B, A\}$ and all points $X$ betweem $A$ and $B$, that is $\mathcal{B}(A, X, B)$.

    This intersection consists of the same points contained in the segment $AB$, as
    \[
    AB = \{A, B\} \cup \{X \in \mathbf{P} \mid \mathcal{B}(A, X, B)\}.
    \]
    Thus, the union of the rays $\overleftarrow{AB}$ and $\overrightarrow{AB}$ is equal to the segment $AB$.
\end{proof}

\subsection{Midpoints and Bisectors}

The \textit{midpoint} of a segment $AB$ is the point $M$ such that $AM \cong MB$ and $[AMB]$.

\begin{axiom}[Midpoint Existence]
    $\forall A,B\in\mathbf{P}, A\neq B \implies \exists M\in\mathbf{P}\colon (\mathcal{B}(A,M,B)) \land (AM \cong MB)$.
\end{axiom}

\begin{axiom}[Midpoint Uniqueness]
    $\forall A,B\in\mathbf{P}, A\neq B \implies \forall M,N\in\mathbf{P}\colon (AM \cong MB) \land (\mathcal{B}(A,M,B)) \land (AN \cong NB) \land (\mathcal{B}(A,N,B)) \implies M = N.$
\end{axiom}

https://chat.deepseek.com/a/chat/s/c1efa70e-bf74-4e92-b967-4054f3472f26

By removing betweenness as a criteria, we obtain the set of points $\{ X \mid AM \cong MB \}$. This set is a line that is perpendicular to the segment $AB$. Alternatively, we can define the perpendicular bisector as the line containing the midpoint $M$ of segment $AB$

% lemma X pg. 4

\section{Parallelism}
Two lines are parallel if they have no point in common. In our set-theoretic foundations, this is to mean that the intersection of the lines is the empty set.
Because we define parallelism as a binary relation on lines, we first must have a set of all lines to work with. Thus, we let $\mathbf{P}$ be the domain called the set of all points and define $\mathbf{L}$ to be the collection of all subsets of $\mathbf{P}$ that consist of the line $\overleftarrow{AB}$ for a any given two points $A, B \in \mathbf{P}$. Formally, that is:
\[
\mathbf{L} \coloneqq \{l \subseteq \mathbf{P} \mid \exists A, B \in \mathbf{P} (A \neq B \land l = \overleftrightarrow{AB})\}.
\]
\begin{definition}[Parallel]
    Given any lines $l, m \in \mathbf{L}$, we say $l$ is parallel to $m$ and write $l \parallel m$ if and only if they are the same same set of points or their set intersection is empty.
    \[
    \forall l, m \in \mathbf{L}, \, l \parallel m \iff (l = m) \lor (l \cap m = \emptyset).
    \]
\end{definition}

add more here

% Define a modern, clean UI color palette (Tailwind CSS inspired)
\definecolor{UIBlueMain}{HTML}{2563EB}    % Blue 600
\definecolor{UIBlueLight}{HTML}{EFF6FF}   % Blue 50
\definecolor{UIBlueBorder}{HTML}{BFDBFE}  % Blue 200
\definecolor{UITextDark}{HTML}{1F2937}    % Gray 800
\definecolor{UITextMuted}{HTML}{6B7280}   % Gray 500
\definecolor{UIAccentRose}{HTML}{F43F5E}  % Rose 500

\begin{tikzpicture}[
    font=\sffamily,
    >={Latex[length=2.5mm, width=2.5mm]},
    % UI Component Styles
    panel/.style={fill=UIBlueLight!50, draw=UIBlueBorder, rounded corners=6pt, thick},
    title/.style={font=\bfseries\color{UITextDark}, anchor=north west},
    subtitle/.style={font=\scriptsize\color{UITextMuted}, anchor=north west},
    % Geometry Element Styles
    point/.style={circle, fill=UIBlueMain, draw=white, thick, inner sep=1.8pt},
    line/.style={thick, UIBlueMain},
    accentline/.style={thick, UIAccentRose}
]

    % ==========================================
    % PANEL 1: Segment
    % ==========================================
    \begin{scope}[shift={(0,0)}]
        \filldraw[panel] (0,0) rectangle (6,3);
        \node[title] at (0.2, 2.8) {Segment $AB$};
        \node[subtitle] at (0.2, 2.3) {$\{A, B\} \cup \{X \in \mathbf{P} \mid \mathcal{B}(A, X, B)\}$};

        \coordinate (A) at (1.5, 1);
        \coordinate (B) at (4.5, 1);

        \draw[line] (A) -- (B);
        \node[point, label={[UITextDark]below:{$A$}}] at (A) {};
        \node[point, label={[UITextDark]below:{$B$}}] at (B) {};
    \end{scope}

    % ==========================================
    % PANEL 2: Ray
    % ==========================================
    \begin{scope}[shift={(6.5,0)}]
        \filldraw[panel] (0,0) rectangle (6,3);
        \node[title] at (0.2, 2.8) {Ray $\overrightarrow{AB}$};
        \node[subtitle] at (0.2, 2.3) {Union of $AB$ and points $X$ where $\mathcal{B}(A, B, X)$};

        \coordinate (A) at (1.5, 1);
        \coordinate (B) at (3.5, 1);
        \coordinate (End) at (5, 1);

        \draw[line, ->] (A) -- (End);
        \node[point, label={[UITextDark]below:{$A$}}] at (A) {};
        \node[point, label={[UITextDark]below:{$B$}}] at (B) {};
    \end{scope}

    % ==========================================
    % PANEL 3: Line
    % ==========================================
    \begin{scope}[shift={(0,-3.5)}]
        \filldraw[panel] (0,0) rectangle (6,3);
        \node[title] at (0.2, 2.8) {Line $\overleftrightarrow{AB}$};
        \node[subtitle] at (0.2, 2.3) {Extends infinitely in both directions};

        \coordinate (Start) at (1, 1);
        \coordinate (A) at (2, 1);
        \coordinate (B) at (4, 1);
        \coordinate (End) at (5, 1);

        \draw[line, <->] (Start) -- (End);
        \node[point, label={[UITextDark]below:{$A$}}] at (A) {};
        \node[point, label={[UITextDark]below:{$B$}}] at (B) {};
    \end{scope}

    % ==========================================
    % PANEL 4: Opposite Rays (Intersection)
    % ==========================================
    \begin{scope}[shift={(6.5,-3.5)}]
        \filldraw[panel] (0,0) rectangle (6,3);
        \node[title] at (0.2, 2.8) {Intersection of Opposite Rays};
        \node[subtitle] at (0.2, 2.3) {$\overleftarrow{AB} \cap \overrightarrow{AB} = AB$};

        \coordinate (Start) at (1, 1);
        \coordinate (A) at (2.5, 1);
        \coordinate (B) at (4.5, 1);
        \coordinate (End) at (5, 1);

        % Ray BA (shifted slightly up)
        \draw[accentline, <-, transform canvas={yshift=3pt}]
            (Start) -- (B) node[midway, above=1pt, font=\scriptsize, UIAccentRose] {$\overleftarrow{AB}$};

        % Ray AB (shifted slightly down)
        \draw[line, ->, transform canvas={yshift=-3pt}]
            (A) -- (End) node[midway, below=1pt, font=\scriptsize, UIBlueMain] {$\overrightarrow{AB}$};

        \node[point, fill=UITextDark, label={[UITextDark]below:{$A$}}] at (A) {};
        \node[point, fill=UITextDark, label={[UITextDark]below:{$B$}}] at (B) {};
    \end{scope}

    % ==========================================
    % PANEL 5: Midpoint & Bisector
    % ==========================================
    \begin{scope}[shift={(0,-7)}]
        \filldraw[panel] (0,0) rectangle (6,3);
        \node[title] at (0.2, 2.8) {Midpoint \& Bisector};
        \node[subtitle] at (0.2, 2.3) {$AM \cong MB$ and perpendicular line $l \perp AB$};

        \coordinate (A) at (1.5, 1);
        \coordinate (B) at (4.5, 1);
        \coordinate (M) at (3, 1);
        \coordinate (T) at (3, 2.2);
        \coordinate (Bot) at (3, 0.2);

        \draw[line] (A) -- (B);
        \draw[accentline, <->] (Bot) -- (T) node[right, font=\scriptsize, UIAccentRose] {$l$};

        % Right angle marker
        \draw[accentline] (3, 1.2) -- (3.2, 1.2) -- (3.2, 1);

        % Congruence Tick marks
        \draw[line] (2.25, 0.85) -- (2.25, 1.15);
        \draw[line] (3.75, 0.85) -- (3.75, 1.15);

        \node[point, label={[UITextDark]below:{$A$}}] at (A) {};
        \node[point, label={[UITextDark]below:{$B$}}] at (B) {};
        \node[point, fill=UIAccentRose, label={[UIAccentRose]below left:{$M$}}] at (M) {};
    \end{scope}

    % ==========================================
    % PANEL 6: Parallelism
    % ==========================================
    \begin{scope}[shift={(6.5,-7)}]
        \filldraw[panel] (0,0) rectangle (6,3);
        \node[title] at (0.2, 2.8) {Parallelism ($l \parallel m$)};
        \node[subtitle] at (0.2, 2.3) {No points in common: $l \cap m = \emptyset$};

        \coordinate (L1) at (1, 1.5);
        \coordinate (R1) at (5, 1.5);
        \coordinate (L2) at (1, 0.7);
        \coordinate (R2) at (5, 0.7);

        \draw[line, <->] (L1) -- (R1) node[right, font=\scriptsize\bfseries, UIBlueMain] {$l$};
        \draw[accentline, <->] (L2) -- (R2) node[right, font=\scriptsize\bfseries, UIAccentRose] {$m$};
    \end{scope}

\end{tikzpicture}

\end{document}
